% Diagramme de prédominance d'un couple acide-base (L'essentiel, chapitre 9)
\begin{tikzpicture}[x=1cm,y=1cm]
\draw[otMuted,line width=0.7pt,-{Stealth[length=2.2mm]}] (0,0) -- (10.6,0) node[pcLegende,anchor=north west,yshift=-1mm] {pH};
\draw[otNavy,line width=0.9pt] (5.2,-0.14) -- (5.2,0.14);
\node[pcLegende,anchor=north] at (5.2,-0.2) {$\pKa$};
\node[pcLegendeG,text=otNavy] at (2.5,0.52) {AH prédomine};
\node[pcLegende,text=otMuted] at (2.5,-0.62) {$[\ce{AH}]>[\ce{A-}]$};
\node[pcLegendeG,text=otGreenDark] at (7.9,0.52) {A$^-$ prédomine};
\node[pcLegende,text=otMuted] at (7.9,-0.62) {$[\ce{A-}]>[\ce{AH}]$};
\draw[otOrange,line width=0.6pt,dash pattern=on 2pt off 1.5pt] (5.2,0.2) -- (5.2,1.05);
\node[pcLegende,anchor=south,text=otOrange!80!black] at (5.2,1.05) {$[\ce{AH}]=[\ce{A-}]$};
\end{tikzpicture}
\par\smallskip\legende{La frontière entre les deux domaines de prédominance est le pKa du couple.}
